\documentclass[a4paper,12pt]{article}
\usepackage[serbian]{babel} % Serbian language support
\usepackage[a4paper, margin=2.5cm]{geometry} % Improved page layout
\usepackage{fancyhdr} % Custom headers/footers
\usepackage{titlesec} % Better section titles
\usepackage{tocloft} % Better TOC spacing
\usepackage{biblatex} % Bibliography management
\usepackage{csquotes} % Recommended for biblatex
\usepackage[T1]{fontenc} % For đ


\setlength{\headheight}{14.5pt} % Fix fancyhdr warning

\pagestyle{fancy}
\fancyhf{} % Clear default header/footer
\fancyhead[L]{\small \leftmark} % Header left
\fancyfoot[C]{\thepage} % Page number in center footer

\titleformat{\section}{\large\bfseries}{\thesection}{1em}{}

\title{Pregled gradiva kursa \\[0.5em] \large Računarstvo i društvo}
\author{Marko Lazarević}
\date{Oktobar 2025}

\renewcommand{\contentsname}{Sadržaj}

\begin{document}

\maketitle
\newpage
\tableofcontents
\newpage

\section{Istorijski pregled razvoja računarstva}
\subsection{Pregled gradiva}

Razvoj računarstva možemo sagledati kroz tri ključne celine:
\begin{itemize}
    \item Razvoj hardvera odnosno fizičkih računara
    \item Razvoj programskih jezika i softvera
    \item Razvoj računarskih mreža i interneta
\end{itemize}

Razvoj ove tri celine se međusobno prepliće i uslovljava.

Razvoj računara tekao je od abakusa i mehaničkih sprava za računanje, Žakardovog tkačkog razboja, preko telegrafa, telefona i radija pa do mašina sa bušenim karticama i elektronskih računara sličnih današnjim.

Razvoj modernih računara krenuo je od vakuumskih cevi, koje su zamenjene tranzistorima. Sa razvojem smanjivala se fizička veličina računara, a povećavale performanse i mogućnosti upotrebe.

Razvoj softvera pratio je razvoj samih računara, a bušene kartice možemo smatrati pretečom modernih programskih jezika. Prvi moderni programski jezici su Fortran, Lisp, COBOL i Algol.

Razvoj računarskih mreža bio je uslovljen stalnom željom za komunikacijom. ARPANET smatramo za preteču današnjeg interneta. 

Današnji razvoj je sve brži i računarstvo ima sve više grana. Neki od najkorišćenijih produkata računarstva u današnje vreme su internet, mobilni uređaji, veštačka inteligencija, blockchain i virtuelna stvarnost.


\subsection{Predložene teme}
\begin{itemize}
    \item Windows vs Linux, zašto je Windows osvojio tržište
    \item Open-source pokret i njegov značaj za razvoj softvera
    \item Razvoj računarskog hardvera: od vakuumskih cevi do kvantnih računara
\end{itemize}

\section{Etika}
\subsection{Pregled gradiva}
Etika je grana filozofije koja se bavi izučavanjem morala. Podoblasti etike su:
\begin{itemize}
    \item Metatetika
    \item Normativna etika
    \item Primenjena etika
\end{itemize}

Etika se bavi pitanjima i problemima u kojima odluka jedne strane može oštetiti drugu stranu, a prvoj doneti korist. 

Teorija etike omogućava pravljenje logičkih argumenata u vezi sa problemom koji se posmatra.  Etičke teorije su:
\begin{itemize}
    \item Subjektivni relativizam - svaki pojedinac sam za sebe odlučuje šta je dobro
    \item Kulturološki relativizam - društvo definiše etičke norme
    \item Teorija duhovnog komandovanja - duhovne knjige i verovanja određuju etičke norme
    \item Etički egoizam - osoba odlučuje tako da dugoročno ostvari maksimalnu dobit
    \item Kantijanizam - ljudi se vode moralnim pravilima koja su jednaka za sve
    \item Utilitarizam postupka - akcija je dobra ako donosi više koristi nego štete
    \item Utilitarizam pravila - prihvataju se pravila koja povećavaju ukupnu sreću svih
    \item Teorija socijalnog ugovora - ljudi prihvataju pravila zarad opšteg dobra, pod uslovom da ih i drugi poštuju
    \item Etika vrline - ispravna akcija je ona koju bi uradila moralna osoba koja poseduje vrline kao što su poštenje, pravednost i lojalnost
\end{itemize}

Razvojem tehnologija otvaraju se nova etička pitanja vezana za same tehnologije i njihovu upotrebu. Etički stavovi ljudi prema tehnici: 
\begin{itemize}
    \item Tehnicizam - stav da je tehnologija korisna kao alat i da će jednog dana rešiti sve tehničke probleme
    \item Optimizam - veruje se da tehnologija ima pozitivan uticaj na ljudsko društvo i da ima moralnu osnovu
    \item Skepticizam - stav da tehnološki napredak donosi štetu ljudskom društvu 
\end{itemize}

Računarska etika se bavi raznim etičkim pitanjima vezanim za upotrebu i razvoj računarstva. Postoje postavljena pravila i smernice za etičko korišćenje računara.

\subsection{Predložene teme}
\begin{itemize}
    \item Etičko hakovanje 
    \item Autorska prava i piraterija
\end{itemize}

\section{Upotreba interneta, spam poruke}
\subsection{Pregled gradiva}

Sve više ljudi ima mobilni telefon ili računar. Samim tim sve više ljudi ima e-mail. Veliki procenat mejlova danas predstavlja nepoželjne mejlove poznatije kao spam. 

Kompanije i pojedinci koriste spam poruke kako bi se jeftinije reklamirali. Drugi sadržaji koji mogu biti u spam porukama su: eksplicitni sadržaj, finansijski sadržaj i razne prevare. Problemi koji se javljaju zbog spam poruka su:
\begin{itemize}
    \item Gubljenje vremena i novca
    \item Zauzimanje memorije
    \item Opterećenje mreže
\end{itemize}

Spameri prikupljaju adrese na razne načine kako bi povećali svoju bazu. Adrese se mogu prikupljati pomoću virusa koji kradu kontakte, kupiti od drugih kompanija ili generisati nasumično. Spameri mogu instalirati botove za spamovanje na nezaštićene servere i tako dodatno sakriti sopstveni identitet. 

Borba protiv spam poruka vodi se pomoću spam filtera. Spam filteri mogu biti postavljeni na nivou internet provajdera, e-mail servisa ili na samom e-mail klijentu korisnika. 

World Wide Web (WWW) nije isto što i internet. Internet je mreža povezanih računara, a WWW je skup povezanih dokumenata na tim računarima. Svaki dokument ima jedinstvenu adresu. Internet se koristi u razne svrhe, kao što su slanje poruka i komunikacija, kupovina, društvene mreže, blogovi, deljenje sadržaja, učenje i edukaciju, video igre, online kockanje, administraciju, kontrolu uređaja. 

\subsection{Predložene teme}
\begin{itemize}
    \item Phishing
    \item Dark web i anonimnost na internetu
\end{itemize}

\section{Cenzura i sloboda govora, deca i neprikladan sadržaj}
\subsection{Pregled gradiva}

Cenzura je pokušaj da se reguliše ili spreči pristup informacijama koje se smatraju štetnim, opasnim ili uvredljivim. Cenzuru su kroz istoriju najviše vršile vlade i religijske institucije. Cenzura može biti direktna, gde druga strana (npr. vlada) ima monopol nad informacijama, nameće licenciranje i registraciju ili ima direktan pregled publikacije. Drugi oblik cenzure je samocenzura, ona predstavlja situaciju kada strana koja poseduje informaciju samu sebe cenzuriše, odnosno sama odlučuje šta će od informacija podeliti a šta ne. 

Internet znatno otežava cenzuru (društvene mreže, Reddit, privatni sajtovi...) jer korisnici često mogu ostati anonimni, a time izbegavaju bilo kakve posledice zbog deljenja informacija. 

Kroz istoriju su se ljudi borili protiv cenzure. Kant je smatrao da aristokratija i Crkva drže monopol nad razmišljanjem. Mil je smatrao da svima treba dati pravo da slobodno govore i iznose svoje mišljenje, iako ono može sadržati neistine. U današnjem društvu sve je teže slobodno govoriti i misliti. Sjajan primer su studenti u Srbiji protiv kojih se vodi ogromna represija zbog želje za pravdom i borbe protiv režima. 

Neprikladan sadržaj i pornografija lako su dostupni na internetu. Za zaštitu na internetu najčešće se koriste veb filteri. To su programi koji sprečavaju određene veb sajtove da budu prikazani.  Filtere najčešće koriste roditelji, kompanije i obrazovne institucije. Drugi problem za decu na internetu je poznat kao Sexting. Sexting označava slanje poruka neprikladnog sadržaja ili fotografija sa malo ili bez odeće. Može se sprečiti zabranom korišćenja određenih aplikacija maloletnim licima ili automatsko uništavanje takvih poruka. 

\subsection{Predložene teme}
\begin{itemize}
    \item Lažni digitalni identiteti i deepfake
    \item Cenzura medija u Srbiji
    \item Uloga društvenih mreža u širenju i kontroli informacija
    \item Automatizovana propaganda: SNS botovi i sloboda govora na internetu
\end{itemize}

\section{Internet prevare i zavisnost od interneta}
\subsection{Pregled gradiva}

Računari i računarske mreže mogu se zloupotrebiti za internet prevare i lažiranje i fabrikovanje podataka. Najčešće internet prevare:
\begin{itemize}
    \item Krađa identiteta - korišćenjem tuđih podataka ostvaruju se prava te osobe
    \item Pecanje (phishing) - slanje poruka u kojima se hitno traži da korisnik unese neke podatke ili uradi neku akciju, a u stvari iza toga stoji prevara
    \item Nigerijska prevara - od žrtve se traži pomoć za transfer velike količine novca uz obećanje da će žrtva dobiti procenat, uz uslov da pošalje podatke svog bankovnog računa
    \item Lažne ocene proizvoda i usluga - ostavljanje lažnih recenzija i ocena kako bi se povećala prodaja proizvoda ili usluge, gde ocene ne odražavaju pravu sliku kvaliteta
    \item Netačne informacije - internet sadrži milijarde veb stranica, ne postoji stroga provera ili recenzija svake od njih pa se često nailazi na stranice sa lažnim i netačnim informacijama
\end{itemize}

Zavisnost od interneta je stvarna bolest današnjice koja se javlja kod sve većeg broja ljudi. Istraživanja su pokazala da se za svaku pristiglu poruku povećava lučenje dopamina (hormon sreće). Simptomi zavisnosti od interneta:
\begin{itemize}
    \item Konstantno razmišljanje o korišćenju interneta
    \item Ostajanje na mreži duže od planiranog
    \item Bezuspešni pokušaji kontrole ili prekida upotrebe interneta
    \item Provođenje mnogo vremena na internetu radi osećaja zadovoljstva
\end{itemize}
Neki od problema koji se javljaju kao posledica zavisnosti od interneta su: glavobolje, problemi sa vidom, depresija, anksioznost, gubljenje osećaja za vreme... 


\subsection{Predložene teme}
\begin{itemize}
    \item Rad od kuće i prevare na internetu: velike zarade za kratko vreme i piramidalne šeme
    \item Baby sensory videos – jeftina zabava ili psihološko usmeravanje ka zavisnosti
\end{itemize}

\section{Bezbednost interneta}
\subsection{Pregled gradiva}

Hakovanje ima dugu istoriju, stariju od savremenog računastva, ali se u današnje vreme najčešće vezuje za dobro poznavanje računarskih sistema i mreže. Haker je osoba koja neovlašćeno pristupa računarima i računarskim mrežama. 

Hakovanje lozinki najčešće se radi metodom grube sile, za kraće lozinke, i metodom rečnika, za duže lozinke. Zbog toga je preporuka da se koristi što više karaktera, da se ne koriste imena i datumi, kao ni standardne reči. Takođe, veoma je poželjno koristiti višefaktorsku autentifikaciju. 

Hakovanje u okviru "besplatne" Wi-Fi mreže je metoda u kojoj haker "uhvati" kolačić koji veb stranica šalje korisniku. Kasnije taj kolačić koristi kako bi imao ista prava i predstavljao se kao taj korisnik. 

Zakonske regulative predviđaju novčane kazne i kazne zatvora za hakere. 

Zlonamerni softver je softver namenjen da nanese štetu krajnjem korisniku. Napadi zlonamernim softverom mogu biti usmereni na: usporavanje računara, uništavanje ili krađu fajlova, preuzimanje kontrole, krađu identiteta...

Načini distribuiranja zlonamernog softvera:
\begin{itemize}
    \item Virusi - programi koji se samorepliciraju i deo su drugog programa koji se naziva domaćin
    \item Crvi - samostalni programi koji se šire na mreži korišćenjem rupa u bezbednosti
    \item Cross-site scripting - napad pri kojem se zlonamerna skripta ubacuje na stranicu i krade podatke korisnika 
    \item Usputno preuzimanje podataka - pri poseti veb stranice se na računar posetioca kopira zlonamerni softver
    \item Trojanski konj - zlonamerni softver je skriven unutar naizgled korisnog programa
    \item Spyware - softver koji upada u sistem i krade poverljive podatke 
    \item Mreža botova - kolekcija međusobno zaraženih uređaja koji se kontrolišu istim zlonamernim softverom
\end{itemize}

Da bi se zaštitili od zlonamernog softvera potrebno je da redovno ažuriramo sistemski i ostali softver kako bi imali najnovije ispravke. Možemo koristiti softver za odbranu od zlonamernog softvera, vatreni zid (firewall) i slično. Veoma je važno da kao korisnici budemo obazrivi i pazimo na svoje aktivnosti na internetu.

\subsection{Predložene teme}
\begin{itemize}
    \item Blue Whale izazov: primer socijalnog inženjeringa i psihološke manipulacije na internetu
    \item Uticaj zero-day ranjivosti na sajber bezbednost
\end{itemize}

\section{Privatnost na internetu}
\subsection{Pregled gradiva}

Privatnost je veoma širok i neodređen pojam, zbog čega ga je teško definisati. Privatnost možemo definisati kao mogućnost pojedinca da ima određen nivo kontrole nad time ko i u kojoj meri ima pristup njegovim podacima i njegovom telu. 

Ostavljanje informacija na internetu može biti voljno i nevoljno. Društvene mreže imaju najveći uticaj na privatnost na internetu. Korisnici sve više kreiraju privatne profile i uklanjaju komentare drugih i svoje ime sa slika na kojima su označeni. Istraživanja pokazuju da uprkos tome koliko neko odluči da deli informacija, odluke njegovih prijatelja takođe imaju značajnu ulogu u privatnosti. 

Kolačići su male tekstualne datoteke koje čuvaju podatke o korisnicima veb sajtova (npr. podatke o prijavi, sadržaju korpe, jezik i slično). Postoje dve vrste kolačića:
\begin{itemize}
    \item Sesijski - kratkotrajni kolačići, brišu se gašenjem pregledača
    \item Trajni - imaju rok trajanja i čuvaju se na računaru
    \begin{itemize}
        \item kolačići prve strane – koriste ih samo veb sajtovi koji su ih napravili
        \item kolačići treće strane – koriste ih sajtovi koji ih nisu napravili, najčešće za praćenje informacija o posetama na drugim stranicama
    \end{itemize}
\end{itemize}
Upotreba kolačića je zakonski regulisana, pa je svaki veb sajt u obavezi da obavesti korisnika o prikupljanju informacija i da dobije pristanak samog korisnika. Korisnici uglavnom ne čitaju politiku privatnosti, niti znaju kako je njihova privatnost zaštićena. Kolačići se mogu koristiti za reklame, praćenje istorije kupovine, istorije lokacije, informacija o uređajima i klikova na linkove...

\subsection{Predložene teme}
\begin{itemize}
    \item Digitalni tragovi i posledice po privatnost korisnika
    \item Praćenje lokacije putem mobilnih uređaja – prednosti i rizici
\end{itemize}


\end{document}
